% Mytheus — one-page project abstract
% Compiles with pdfLaTeX on Overleaf. No external assets, no custom classes.

\documentclass[10pt,a4paper]{article}

\usepackage[margin=1.8cm]{geometry}
\usepackage{mathptmx}          % Times text + math
\usepackage[T1]{fontenc}
\usepackage[utf8]{inputenc}
\usepackage{microtype}
\usepackage{titlesec}

% Compact, unnumbered headings so the whole piece stays on one page.
\titleformat{\section}{\normalfont\bfseries\normalsize}{}{0pt}{}
\titlespacing*{\section}{0pt}{0.8\baselineskip}{0.2\baselineskip}

\setlength{\parskip}{0.28\baselineskip}
\setlength{\parindent}{0pt}
\pagestyle{empty}

\begin{document}

\begin{center}
  {\LARGE\bfseries Mytheus}\\[0.35em]
  {\large Personalized Multi-Agent Meeting Orchestration}
\end{center}

\vspace{0.4em}
\hrule
\vspace{0.5em}

\section{The problem}

Meeting-assistance tools produce one shared artefact, a transcript and a
summary, and hand the identical copy to everyone. That is the wrong unit of
output: what a designer needs from an hour-long planning call is not what the
engineering lead needs. The tools are also passive, reporting after the fact,
so a decision that contradicts last month's decision, a task assigned to
someone already at capacity, or a reference to context a participant does not
have all pass unremarked while the meeting is still running and still fixable.

\section{What the system does}

A meeting runs as a LiveKit video room whose audio is streamed server-side to
Deepgram, which returns diarized transcripts; interim results are mutable and
only finals are persisted. Live captions render beside the video within two
seconds. Next to them, each participant sees a sidebar that is theirs alone:
notes filtered to their role, decision and action cards awaiting confirmation,
suggested questions ranked by priority, and proactive warnings, with a
red conflict banner pinned above everything else. Confirmed action items
synchronize outward to the team's existing tools, with the connected surface
deliberately capped at Google Calendar, Slack, and one of Jira or Linear.

\section{The agent layer}

Each participant gets their own AI agent. Six agent types share one LangGraph
\texttt{BaseAgent} lifecycle and a versioned prompt loader.
\textbf{PersonalListenerAgent} scores each window of transcript for relevance
to one user's role and skills and records only what matters to them.
\textbf{DecisionAgent} and \textbf{ActionAgent} extract decisions and tasks with
rationale, assignee, and a confidence score. \textbf{QuestionAgent} proposes
role-specific questions while the topic is still open.
\textbf{ProactiveAgent} detects contradictions with past decisions, workload
overload, and missing stakeholders.

The distinctive component is \textbf{ContextGapAgent}. When the conversation
depends on context a participant lacks, that participant's agent queries
\emph{other participants' agents} over a correlation-ID request/response
protocol, falling through a fixed chain: peer agents, then agents from related
meetings, then retrieval, then silence. Agents collaborate across meetings to
close gaps no single attendee could close alone.

Two rules constrain every agent. Confidence thresholds are uniform: above $0.7$
surface to the user, between $0.4$ and $0.7$ log silently, below $0.4$ discard.
And there are \textbf{no silent commits}. An agent emits a \emph{candidate}
only; the database row, the event, and the knowledge-graph edge are written
solely on user confirmation, so nothing is recorded that the user did not
approve.

\section{Organizational memory}

Every confirmed decision and action becomes an edge in a property graph
(Apache AGE) whose nodes are Person, Meeting, Decision, Action, and Topic,
linked by \texttt{ATTENDED}, \texttt{PRODUCED}, \texttt{ASSIGNED\_TO},
\texttt{RELATES\_TO}, and \texttt{FOLLOWS\_FROM}. The graph is what makes
cross-meeting reasoning possible: contradiction detection is a graph traversal,
not a prompt. In parallel, transcript and document chunks are embedded with
BGE-small-en into per-tenant vector collections. Retrieval runs a fixed
five-stage pipeline, expand, retrieve, rerank, assemble, cite: hybrid vector
and full-text results are fused, reordered by a cross-encoder reranker, and
every returned passage carries a citation back to its source chunk.

\section{System architecture}

The platform is event-driven microservices, cut into families that scale on
different signals: a CRUD control plane (requests per second), a real-time
plane (concurrent connections), and an AI/inference plane (tokens and vectors),
with the embedding and reranking models served on GPU off the API tier.
Services communicate over tenant-partitioned Redpanda topics with Avro schemas,
at-least-once delivery, transactional outboxes, and consumer-side
deduplication. A separate real-time gateway fans captions and agent output to
clients over WebSockets, independent of the media path.

Multi-tenancy is enforced structurally rather than in application code. One
chokepoint resolves \texttt{org\_id} from the verified session token and
injects it server-side; PostgreSQL row-level security is the fail-closed
backstop; vector collections are per tenant and the graph is org-partitioned.
No client ever derives or transmits a tenant identifier.

\section{Current status}

The contract scaffold and the foundation-and-tenancy band are built across
eight repositories: the row-level-security baseline with least-privilege
database roles, role-based access grants behind a deny-by-default authorization
gate, append-only audit logging, meetings CRUD, Clerk identity with webhook
user sync, the transactional outbox relay, and both clients reading the real
backend contract. Meeting capture and transcription come next, then the agent
runtime and retrieval.

\end{document}
